Let $f\in\mathbb{C}[z]$ be monic, non-constant, of degree $d\ge 1$. Define the (closed) lemniscate
\[
E(f):=\{z\in\mathbb{C}: |f(z)|\le 1\}.
\]
Interpretation issue: the phrase ``covered by circles'' is ambiguous.

\begin{itemize}
\item \emph{Literal (1-dimensional) interpretation:} circles mean circumferences. Then $E(f)$ typically has nonempty interior (e.g. for $f(z)=z$), so it cannot be covered by countably many circles of finite total radius. Hence that literal reading is impossible.
\item \emph{Standard complex-analysis interpretation (Cartan/Pommerenke context):} ``circles'' means closed discs (each disc determined by its boundary circle). We adopt this corrected interpretation: we ask whether there exist closed discs $\overline{D}(c_j,r_j)$ such that
\[
E(f)\subseteq \bigcup_j \overline{D}(c_j,r_j)
\qquad\text{and}\qquad \sum_j r_j\le 2.
\]
The covering family is understood to be finite or countable so that $\sum_j r_j$ is defined.
\end{itemize}

Edge cases: $d=1$ gives a single closed unit disc, covered with total radius $1$. For $d=2$ the trivial cover by unit discs about the roots already has total radius $2$.
